\section{Implication is a function type}\label{sec:03-propositions-and-proofs:04-implication:1}

\(P \to Q\) (read ``\(P\) implies \(Q\)'') is literally a function type: a proof of \(P \to Q\) is a function that turns any proof of \(P\) into a proof of \(Q\). This is the same \texttt{→} flagged back in \hyperref[sec:01-basics:02-def-let-implicit:1]{Chapter 1, Section 2}, the one that built \texttt{double : Nat → Nat}. Nothing about the symbol changes here, only what stands on either side of it. \texttt{Nat → Nat} is functions between two data types; \texttt{P → Q}, for \texttt{P Q : Prop}, is functions between two \emph{proofs}, which is exactly what ``implication'' turns out to mean.

\begin{lstlisting}
theorem modus_ponens {P Q : Prop} (hpq : P (*$\rightarrow$*) Q) (hp : P) : Q :=
  hpq hp
\end{lstlisting}

\begin{mathreading}

Under Curry--Howard, the implication \(P \Rightarrow Q\) \emph{is} the function space \(P \to Q\) (the set of proofs of \(Q\) parameterized by proofs of \(P\)). So modus ponens \[
\frac{P \Rightarrow Q \qquad P}{Q}
\] is nothing but function application: given \(f \in \mathrm{Hom}(P, Q)\) and \(p \in P\), evaluate to get \(f(p) \in Q\). The term \texttt{hpq hp} is precisely this evaluation \(f(p)\).

\end{mathreading}

\hyperref[sec:03-propositions-and-proofs:03-theorem-lemma:1]{← \texttt{theorem}/\texttt{lemma}} \textbar{} \hyperref[sec:03-propositions-and-proofs:00-index:1]{Index} \textbar{} \hyperref[sec:03-propositions-and-proofs:05-and-or-not:1]{Next: And, Or, Not →}
